\begin{ExerciseBlock}
\ExerciseHeading{}{Verständnisfragen}{free_oscillations_checks}
Entscheiden Sie, welche Aussagen zutreffen, und begründen Sie Ihre Entscheidung.
\begin{enumerate}
\item Harmonische Schwingungen mit größerer Amplitude haben eine längere Periodendauer.
\item Kreisfrequenz und Frequenz sind proportional.
\item Nur ein Federpendel kann harmonisch schwingen.
\item Die Frequenz einer unterkritisch gedämpften harmonischen Schwingung ist kleiner als die des ungedämpften Systems.
\item Bei der Herleitung für ein vertikales Federpendel wurde die Gewichtskraft vergessen.
\item Der Betrag der Beschleunigung ist beim Durchgang durch die Ruhelage am größten.
\item Bei gleicher Federkonstante nimmt die Periodendauer mit der Masse zu.
\item Der Faktor $1/2$ in der Federenergie entsteht, weil die Federkraft mit der Auslenkung linear ansteigt.
\item Ein Motorboot wird hauptsächlich durch Stokes-Reibung gebremst.
\item Je stärker die Dämpfung ist, desto schneller kehrt das System in die Ruhelage zurück.
\item Ein mathematisches Pendel schwingt auch bei großen Auslenkungen harmonisch.
\end{enumerate}
\end{ExerciseBlock}
\begin{Solution}
\par\medskip\noindent\textbf{Musterlösung}\par\nopagebreak[4]
\begin{enumerate}
\item Falsch. Die Periode eines linearen Oszillators ist unabhängig von der Amplitude.
\item Wahr. Es gilt $\omega=2\pi f$.
\item Falsch. Auch ein Drehpendel oder ein Fadenpendel bei kleinen Winkeln kann harmonisch schwingen.
\item Wahr. Im Schwingfall gilt $\omega=\sqrt{\omega_0^2-\delta^2}$.
\item Falsch. Die Gewichtskraft verschiebt die statische Ruhelage. Die Bewegungsgleichung für die Auslenkung relativ zu dieser Lage enthält sie bereits.
\item Falsch. Wegen $a=-\omega_0^2x$ ist die Beschleunigung in der Ruhelage null.
\item Wahr. Es gilt $T_0=2\pi\sqrt{m/D}$.
\item Wahr. Integration der Rückstellkraft im Betrag ergibt $\int_0^x D\xi\,\mathrm d\xi=Dx^2/2$.
\item Falsch. Die Stokes-Näherung setzt kleine Reynolds-Zahlen voraus.
\item Falsch. Bei starker Überdämpfung kann die Rückkehr deutlich langsamer werden.
\item Falsch. Das Rückstellmoment ist proportional zu $\sin\varphi$. Erst die Kleinwinkelnäherung macht es linear.
\end{enumerate}
\end{Solution}
